\documentclass[12pt]{book}
\usepackage[utf8]{inputenc}
\usepackage[T1]{fontenc}
\usepackage{amsmath,amssymb,amsfonts}
\usepackage{amsthm}

\begin{document}

\setcounter{page}{381}
We deduce the following elementary but essential consequences:
(i)If \(\theta_{f}\) and \(\theta_{g}\) are isomorphisms for all arguments, so is \(\theta_{gf}\).
(ii)If \(\theta_{g}\) and \(\theta_{gf}\) are isomorphisms for all arguments, so is \(\theta_{f}\).
(iii)If \(\theta_{f}\) and \(\theta_{gf}\) are isomorphisms for all arguments, then \(\theta_{g}\) is an isomorphism for every complex of the form \(\mathbf{R} f_{*} F^{*}\) with \(F^{*} \in D_{c}^{-}(X)\).

e) By noetherian induction on \(X\), we may assume the theorem proven for every morphism of the form \(g: Z \to Y\), where \(g=fi\), and \(i: Z \to X\) is a closed immersion of \(Z\) onto a subscheme of \(X\), different from \(X\).

f) We now apply Chow's Lemma [EGA II 5.6.1] to deduce the existence of a scheme \(X'\) projective over \(Y\), together with a projective \(Y\)-morphism \(g: X' \to X\), which is an isomorphism on a non-empty open subset \(U\) of \(X\) (we may assume \(X \ne \emptyset\)).

Given a coherent sheaf \(F\) on \(X\), consider the natural map

\setcounter{page}{382}
\[F \to \mathbf{R}g_{*} g^{*} F,\]
and embed it in a distinguished triangle.
Then since \(g\) is an isomorphism on \(U\), \(G^{*}\) has support on \(X-U\), and so \(\theta_{f}(G^{*})\) is an isomorphism by our induction hypothesis e).
Thus it will be sufficient to show that \(\theta_{f}(\mathbf{R g}_{*} g^{*} F)\) is an isomorphism.
Using (iii) above, we reduce to showing \(\theta_{g}\) and \(\theta_{fg}\) are isomorphisms for all arguments, and so we reduce to the case \(f\) is projective.

g) Since \(Y\) is affine, we can embed \(X\) in a suitable projective space over \(Y\), say \(f=pi\), with
\(p: \mathbb{P}_{Y}^{n} \to Y\) the projection.
Using (i) above, we can treat \(i\) and \(p\) separately.
Now \(i\) is a finite morphism, so \(\theta_{i}\) is an isomorphism by Proposition 3.1 above, and [III 6.7].
Also \(\theta_{p}\) is an isomorphism by Proposition 3.2 above, and [III 5.1].
q.e.d.

\setcounter{page}{383}
We can now pull ourselves up by our bootstraps, and obtain a theory of \(f^{!}\) and \(Tr_{f}\) for complexes with coherent cohomology and schemes admitting dualizing complexes.

Corollary 3.4. We consider the category of noetherian preschemes which admit a dualizing complex, and we consider morphisms of finite type. Then

(a). For every such morphism \(f: X \to Y\), there is a theory of
functor
\[f^{!}: D_{c}^{+}(Y) \to D_{c}^{+}(X)\]
plus the data 2)-5) and properties VAR 1 - VAR 6 of [III 8.7] (only leave out the word "embeddable" wherever it occurs).

(b). For every such proper morphism \(f: X \to Y\), there is a theory of trace consisting of a functorial morphism
\[Tr_{f}: \mathbf{R} f_{*} f^{!} \to 1\]
with the properties TRA 1 - TRA 4 of [III 10.5] (only leave out the phrase "projectively embeddable" wherever it occurs).

\setcounter{page}{384}
(c). For every such proper morphism \(f: X \to Y\), the duality morphism
\[\underline{\theta}_{f}: \mathbf{R} f_{*} \underline{Hom}_{X}^{\cdot}\left(F^{*}, f^{!} G^{*}\right) \to \mathbf{R} \underline{Hom}_{Y}^{\cdot}\left(\mathbf{R} f_{*} F^{*}, G^{*}\right),\]
obtained by composing the morphism of [II 5.5] with \(Tr_{f}\), is an isomorphism for \(F^{*} \in D_{qc}^{-}(X)\) and \(G^{*} \in D_{c}^{+}(Y)\).
(Compare [I 11.1].)

Proof. (a)
Let \(K^{*}\) be a residual complex on \(Y\) (and is then a dualizing complex on \(Y\), and \(Q f^{\Delta} K\) a dualizing complex on \(X\)).
Let \(\underline{D}_{Y}\) and \(\underline{D}_{X}\) be as in the theorem above, and define
\[f^{!}\left(G^{*}\right)=\underline{D}_{X}\left(\mathbf{L} f^{*} \underline{D}_{Y}\left(G^{*}\right)\right)\]
for all \(G^{*} \in D_{C}^{+}(Y)\).
Note that \(\underline{D}_{Y}(G^{*}) \in D_{C}^{-}(Y)\), so that \(\mathbf{L} f^{*}\) is defined.
Here is where we need that \(QK^{*}\) is a dualizing complex, not just pointwise dualizing.
Observe also that \(f^{!}\) is independent of the choice of \(K^{\cdot}\) (use [V 3.1]).

The construction of the isomorphisms \(c_{f,g}\) is easy, using [VI 3.1].
For the isomorphisms \(d_{f}\) and \(e_{f}\) we use [III 6.9b] and [III 2.4], respectively.
The details of these constructions, and verifications of VAR 1 - VAR 6 are left to the reader (:).

Observe, by the way, that this part of the Corollary does not depend on the duality theorem, and could have come just after the construction of \(f^{\Delta}\) [VI 3.1].

\setcounter{page}{385}
(b) To define \(Tr_{f}\), let \(G^{\cdot} \in D_{C}^{+}(Y)\), and use [II 5.10],
[VI 4.2], Theorem 2.1 above, and [V 2.1].
The verification of TRA 1 is clear, using TRA 1 of [VI 4.2].
For TRA 2 we use [III 6.9c]. Details left to reader (:).

(c) To prove the duality formula, we reduce as in part b) of the proof of the theorem above, to the case of complexes with coherent cohomology.
Let \(F^{*} \in D_{c}^{-}(X)\) and \(G^{*} \in D_{C}^{+}(Y)\) Consider
\[\mathbf{R} f_{*} \mathbf{Hom}_{X}\left(F^{*}, f^{*} G^{*}\right) . \tag{1}\]

We write \(F^{*}=\underline{D}_{X} \underline{D}_{X}(F^{*})\) and \(f^{\cdot} G^{\cdot}=\underline{D}_{X}(\mathbf{L}^{*} \underline{D}_{Y}(G^{\cdot}))\).
Now since \(\underline{D}_{X}\) is a dualizing functor, it transforms \(\mathbf{R}Hom\) into \(\mathbf{R}Hom\) of the duals of the arguments, with the order reversed.

\setcounter{page}{386}
Thus (1) becomes
\[\mathbf{R} f_{*} \mathbf{Hom}_{X}\left(\mathbf{L} f^{*} \underline{D}_{Y}\left(G^{\cdot}\right), \underline{D}_{X}\left(F^{\cdot}\right)\right) . \tag{2}\]

Applying the duality theorem above, this becomes
\[D_{Y} \mathbf{R} f_{*} \mathbf{Hom}_{X}\left(\mathbf{L} f^{*} \underline{D}_{Y}\left(G^{\cdot}\right), \underline{D}_{X}\left(F^{\cdot}\right)\right) \tag{3}\]
which in turn, since \(\mathbf{L} f^{*}\) is a left adjoint of \(\mathbf{R} f_{*}\),
\[D_{Y} \mathbf{Hom}_{Y}\left(\underline{D}_{Y}\left(G^{\cdot}\right), \mathbf{R} f_{*} \underline{D}_{X}\left(F^{\cdot}\right)\right) . \tag{4}\]

Now applying duality to the second argument, (4) becomes
\[\mathbf{R}Hom_{Y}\left(\mathbf{R} f_{*} F^{*}, G^{*}\right) .\]

The reader may check that our chain of isomorphisms is indeed \(\underline{\theta}_{f}\), which proves (c).

\setcounter{page}{387}
Proposition 3.5 (Compatibility of Local and Global duality). Let \(f: X \to Y\) be a proper morphism of noetherian schemes,
let \(x \in X\) be a closed point, with local ring \(A=\sheaf{O}_{x}\).
Let \(R^{*}\) be the dualizing complex \((Qf^{A} K^{*})_{x}\) on \(A\), and assume that \(R^{*}\) is normalized.
Let \(I=\Gamma_{x}(R^{*})\).
Let \(F^{*} \in D_{c}^{+}(X)\).
Then the diagram
is commutative, where \(\theta_{x}\) is the local duality isomorphism, \(\alpha\) is the natural map of derived functors obtained from the inclusion \(\Gamma_{x} \subseteq f_{*}\),
and the trace map
\[Tr_{f, x}: f_{*} J(x)=I \to K^{*},\]
and \(\underline{D}_{Y}\) is the transpose by \(\underline{D}_{Y}\) of the global duality isomorphism of Theorem 3.3.
(Note we write four times Hom instead of \(\mathbf{R}Hom\), because the second argument in each case is injective.)

Proof. Immediate from the definitions of the maps in question.

Remark. This compatibility is the one needed to complete the proof of "Lichtenbaum's theorem" [LC theorem 6.9, see parenthetical remark in middle of p. 103].

\end{document}